% 本文件为 2026 年课程约束与非功能需求审阅稿。
\begingroup
\color{red}

\subsection{安全与隐私补充}

\textbf{传输安全纠正：}HTTP POST 只是不把参数放在 URL 中，本身不能保证密码机密性。部署和验收环境必须使用 HTTPS；密码仅在服务端使用 BCrypt 等强哈希保存，日志、异常和接口响应不得输出明文密码、完整 Token 或钱包敏感信息。

系统采用 RBAC，对 USER、BUSINESS、RIDER、ADMIN 分别授权。资源访问同时校验角色与资源归属，避免顾客读取他人订单、商家操作非所属店铺、骑手查看非本人任务或管理员接口被普通用户调用。

骑手只在履约期间获得完成任务所需的最小地址和联系方式；任务完成后限制继续读取。用户偏好、位置、语音和图片属于敏感或个人数据，需说明用途并允许用户清除。

钱包、积分、优惠券和订单状态变更必须保存审计流水。后台调整必须记录管理员、原因、变更前后值和时间，禁止直接覆盖余额而不留痕。

\textbf{拦截器纠正：}401 表示认证失效，应清理登录状态或引导重新登录，不应无条件重试；网络或 5xx 错误只对幂等请求做有限退避重试，支付、接单、核销等写操作不得在没有幂等键时自动重试。

\subsection{REST 接口与架构约束}

所有新增接口使用 \texttt{/api/v1} 版本前缀；存量接口在不破坏基础功能的前提下逐步迁移。URI 使用资源名词，GET 只查询，POST 创建资源或表达明确业务动作，PATCH 局部更新，DELETE 执行业务逻辑删除。

允许使用 \texttt{/orders/\{id\}/cancel}、\texttt{/delivery-tasks/\{id\}/accept} 等业务动作接口，但不得使用 GET 完成加入购物车、删除商品、提交订单或修改状态。

统一响应至少包含 \texttt{code}、\texttt{msg}、\texttt{data} 和便于排错的 \texttt{traceId}；同时使用 200/201/204、400、401、403、404、409、422、500 等标准 HTTP 状态码。分页、筛选、排序和时间范围使用查询参数。

后端严格保持 Entity、Mapper/Dao、Service、Controller 四层：Controller 只负责接收、校验和响应；业务状态、优惠计算、权限和事务放在 Service；Mapper 只负责数据访问。地图、AI、语音、识图和对象存储通过适配器接口隔离，便于替换或模拟。

\subsection{性能、可靠性与兼容性}

以下指标作为小组验收候选基线，最终数值需在测试环境确定后确认：

\begin{longtable}{@{}p{0.22\linewidth}p{0.61\linewidth}@{}}
\caption{非功能验收候选指标}\tabularnewline
\toprule
\textbf{质量属性} & \textbf{候选验收标准} \\
\midrule
\endfirsthead
\toprule
\textbf{质量属性} & \textbf{候选验收标准} \\
\midrule
\endhead
\bottomrule
\endfoot
普通接口性能 & 在 50 个并发用户的课程测试环境中，非 AI 普通接口 P95 响应时间不超过 1 秒，错误率低于 1\%。 \\
页面响应 & 店铺、商品、订单等主要页面在稳定网络下 3 秒内显示首屏有效内容，并提供加载、空数据和失败状态。 \\
统计性能 & 在不少于 1 万条订单的测试数据下，单个用户或商家的一年看板查询在 3 秒内返回。 \\
外部服务 & AI、语音、识图和地图调用设置超时；超过 10 秒返回降级结果，不无限等待。 \\
并发一致性 & 并发抢单、库存扣减、优惠核销、钱包扣款和积分发放不出现重复成功、负余额或超卖。 \\
事务可靠性 & 订单、资产、优惠、库存任一关键步骤失败时事务回滚或进入可恢复状态，不产生半完成订单。 \\
兼容性 & 支持当前主流 Chromium 内核桌面浏览器；核心顾客与骑手操作在常见移动端宽度下可完成。 \\
可观测性 & 关键请求携带 traceId；登录、订单、配送、资产和外部服务错误具有分级日志且不泄露敏感数据。 \\
\end{longtable}

\subsection{可用性与可访问性}

顾客、商家、骑手和管理员进入与自身角色匹配的工作台，状态名称在前后端保持一致。关键交易按钮在重复点击时不可造成重复提交，并提供处理中、成功、失败和恢复路径。

主题换肤必须保证文本与背景有足够对比度；订单和配送状态不能只依赖颜色表达。表单错误应定位到具体字段，语音、识图和地图功能均提供键盘或文字替代流程。

\subsection{可测试性与 TDD 约束}

所有接口和核心业务方法必须遵循：分析需求与验收条件、先提交失败测试、实现最小功能、测试全部通过、重构并再次通过测试。禁止先开发后补测试，禁止没有断言的形式化测试。

核心业务接口必须全部有自动化测试，关键业务方法覆盖率不低于 90\%。每个需求至少覆盖正常、异常/权限和边界场景；抢单、库存、钱包、优惠券和积分增加并发及重复请求场景。

测试分层包括：Service 单元测试、Controller/MockMvc 接口测试、Mapper/数据库集成测试、前端组件或流程测试、ApiFox/Postman 交叉验收。外部 AI、地图、语音、识图和 OSS 在自动化测试中使用可控桩或模拟服务。

测试类、覆盖率报告、运行日志、需求变更后的回归结果必须纳入仓库。Git 提交顺序应能证明测试早于对应实现，4 名成员均应有与分工一致的持续提交记录。

\section{需求分期与验收追踪}

\subsection{三周实施边界}

\begin{longtable}{@{}p{0.16\linewidth}p{0.22\linewidth}p{0.45\linewidth}@{}}
\caption{需求分期与交付证据}\tabularnewline
\toprule
\textbf{阶段} & \textbf{需求范围} & \textbf{必须留存的证据} \\
\midrule
\endfirsthead
\toprule
\textbf{阶段} & \textbf{需求范围} & \textbf{必须留存的证据} \\
\midrule
\endhead
\bottomrule
\endfoot
第 1 周：基础固定需求 & 用户、商家、店铺、商品、购物车、地址、基础订单 & 基线 SRS、接口契约、失败测试提交、实现提交、基础测试日志与 100\% 核心接口测试。 \\
第 2 周：需求变更 & 骑手主链路及 P1 创新模块；P2 外部能力按风险推进 & 变更影响分析、新增失败测试、实现提交、旧功能回归、架构无需大面积重构的说明。 \\
第 3 周：完善验收 & 重构、异常校验、交互、性能、文档、交叉验收；补齐可完成的 P2 & 重构前后测试、覆盖率、接口文档、AI 使用记录、交叉验收记录、问题与修复清单。 \\
\end{longtable}

\subsection{需求追踪规则}

每项需求使用稳定 ID，并在追踪矩阵中关联：需求说明、优先级、负责成员、接口/页面、数据表、测试类与测试方法、首次失败测试提交、实现提交、回归结果和验收状态。

当教师发布随机差异化需求时，先新增需求 ID 和验收条件，再补充失败测试；不得直接修改实现后再反向编写说明。状态机、优惠策略、推荐服务和第三方适配器应通过接口或配置扩展，避免大面积重构。

\subsection{四人协作约束}

4 名成员可分别主责以下四个创新工作包，但每个工作包都必须包含需求、数据库与接口、前端、自动化测试和验收证据，不能按前后端简单割裂：

1. 营销资产模块：优惠券叠加、会员充值、钱包、积分、资金与资产流水、退款回退及并发幂等。
2. 个性化体验模块：主题换肤、每日幸运色、主题持久化、口味/忌口/价格偏好及隐私授权。
3. 智能交互模块：智能推荐、图片识菜、语音点单草稿、导航适配、第三方服务超时与降级。
4. 评价与数据看板模块：订单评价、商家回复、内容治理、用户消费分析、商家热销/销量统计和商品详情增强。

骑手配送不是第五个可选创新点，而是新增的公共核心链路。小组需指定一名主责人维护骑手角色、配送任务和订单状态机，其余成员分别完成钱包结算、导航、通知、评价和统计等与骑手链路的接口集成。每位业务开发者必须先编写自己模块的失败测试；模块主责人负责正常流程，交叉评审人补充权限、边界、并发和回归场景。

小组需要记录 AI 辅助使用点，包括生成过的脚手架、测试思路和 Bug 定位建议；核心架构决策、测试断言、验收判断和业务代码解释必须由成员完成。

\subsection{完成定义}

单项需求只有同时满足以下条件才可标记完成：SRS 与接口契约已更新；失败测试先于实现存在；正常、异常和边界测试通过；权限、事务和日志符合要求；前后端联调通过；旧功能回归通过；代码与文档已由另一名成员评审；Git 和测试证据可追溯。

\endgroup
